\section{Точные вычисления для гармонических моделей}
Положим $C_n(t)=2T_n(t/2)$. Все приведённые ниже проверки выполнены
над точными рациональными коэффициентами в SymPy 1.14.0. Непрохождение
ограниченного алгоритма не означает невыпрямляемость.

\subsection{Перебор понижающих элементарных преобразований}
Если $c=ra+sb$ при $r,s\geq1$, то
\[
 (x,y,z)\longmapsto(x,y,C_r(x)C_s(y)-z)
\]
есть полиномиальная инволюция и переводит $(C_a,C_b,C_c)$ в
$(C_a,C_b,C_{|ra-sb|})$. Поскольку $0<|ra-sb|<c$, сумма индексов
строго убывает. Скрипт \texttt{reduction\_scan.py} перебирает \emph{все}
представления $c=ra+sb$ на каждом шаге, сортирует новую тройку и
мемоизирует множество всех достижимых терминальных троек. Поэтому
это полный перебор строго понижающих путей указанного вида, а не
жадный выбор единственного представления.

Тройка с индексом 1 непосредственно выпрямляется над $\mathbb Q$.
Достижение $(3,4,5)$ или $(3,5,7)$ даёт выпрямление над
$K=\mathbb Q(\sqrt{-3})$ при использовании проверенных сертификатов
репозитория. Других преобразований алгоритм не использует.
\begin{center}
\begin{tabular}{r|r|r|r|r}
Максимум & Всего & Достигнут индекс 1 & Базы 345/357 & Не покрыто \\
\hline
20 & 297 & 82 & 13 & 202\\
50 & 5382 & 734 & 72 & 4576\\
100 & 44142 & 3209 & 237 & 40696
\end{tabular}
\end{center}
Здесь перечислены строго возрастающие попарно взаимно простые тройки
$2\leq a<b<c\leq N$. Категории не пересекаются: при достижении
индекса 1 тройка учитывается в первом положительном столбце.
Это нижняя граница числа выпрямляемых моделей, а не классификация
остальных моделей. В частности, алгоритм не исследует пути с временным
увеличением степеней и произвольные полиномиальные автоморфизмы.
Для $a=3$ соответствующие числа при $N=100$ равны
$1409=425+142+842$. Первые непокрытые ядра при $a=3$:
$(3,7,8),(3,7,11),(3,8,13),(3,10,11),(3,10,17)$.
Другой важный первый пример: $(4,5,7)$.

\subsection{Матрица при кубической проекции}
Пусть $x=C_3(t)$, $s=C_2(t)=t^2-2$ и
$U_j=U_j(x/2)$ --- многочлен Чебышёва второго рода, причём $U_{-1}=0$.
Из рекурсии $C_{n+3}=xC_n-C_{n-3}$ следует
\[
 C_{3q+1}=U_qt-U_{q-1}s,\qquad
 C_{3q+2}=-U_{q-1}t+U_qs.
\]
Таким образом, $C_b=A_b(x)t+B_b(x)s$, и любые две координаты
дают явную матрицу $N_{b,c}$ над $\mathbb Q[x]$.
Для первых моделей её определители равны
\[
\begin{array}{c|c}
(b,c)&\det N_{b,c}\\\hline
(4,5)&x^2-1\\
(5,7)&-x(x^2-2)\\
(7,8)&x^4-3x^2+1\\
(7,11)&x(x^2-1)(x^2-3)
\end{array}
\]
Поэтому рациональное восстановление $(t,s)$ элементарно. Обращение
определителя не доказывает существование полиномиального автоморфизма:
его нули являются конкретными исключительными слоями, требующими
проверки.

\subsection{Строгая обструкция одного естественного переноса для $H(3,7,8)$}
Тождество
\[
 C_2(C_3)C_8-C_7^2+2=C_2
\]
даёт модификацию
\[
 \beta(x,y,z)=(x,(x^2-2)z-y^2+2,y),
\]
которая переводит $(C_3,C_7,C_8)$ в $(C_3,C_2,C_7)$.
Последняя тройка допускает относительную координату
$H_0(x,Y,Z)=xY^2-2x-Z$, поскольку
$C_3(C_2^2-2)-C_7=t$. Однако
\[
 H_0(x,2-y^2,y)=x(2-y^2)^2-2x-y
\]
имеет степень 4 по $y$ при $x^2=2$.

Более сильно, замена $H_0$ на другой многочлен $H$, удовлетворяющий
$H(C_3,C_2,C_7)=t$, не может сделать центр линейным на обоих
исключительных слоях. Действительно, пусть $r^2=2$ и предположим
$H(r,2-y^2,y)=\alpha y+\beta$.
В $k[t]/(t^3-3t-r)$ имеем
\[
 C_7=t-r(t^2-2),\qquad 2-C_7^2=C_2.
\]
Следовательно, $t=\alpha(t-rt^2+2r)+\beta$.
Базис $1,t,t^2$ даёт одновременно $\alpha=1$ и $\alpha r=0$,
что невозможно. Это доказывает невозможность применения именно
линейного критерия центра к этой модификации и этой параметризации.
Это \textbf{не} доказательство невыпрямляемости $H(3,7,8)$.

\subsection{Нижняя граница для относительной координаты $H(3,7,8)$}
Все многочлены общей степени не выше 4, удовлетворяющие
$H(C_3,C_7,C_8)=t$, образуют двумерное аффинное семейство
$H=H_0+\lambda F_0+\mu F_1$, где
\begin{align*}
H_0={}&x^3z+x^3-x^2yz-xy^2-2xz-x+y^3+2yz-2y,\\
F_0={}&x^3y-x^2z-3xy+y^2-z^2+2z,\\
F_1={}&x^4-x^2z^2+x^2z-3x^2+xy^3-xyz-2xy\\
 &+y^2z-y^2-z^3+2z^2-z+2.
\end{align*}
Это полное решение конечной линейной системы по всем мономам общей
степени не выше 4. Для любого $\lambda,\mu$ справедливо
\[
 H_y(0,0,1)=H_z(0,0,1)=0,\qquad H_x(0,0,1)=-3.
\]
Поэтому $(x,H)$ нельзя дополнить до полиномиального автоморфизма,
сохраняющего $x$. Следовательно, относительная координата параметра,
если существует, должна иметь общую степень как минимум 5.
Это утверждение касается только автоморфизмов, сохраняющих $x=C_3$;
оно не исключает произвольного выпрямления с параметрической
компонентой степени 4.

\subsection{Ограниченный квадратичный поиск}
Для координат $(x,y,z)=(C_3,C_8,C_7)$ положим
$k=x^2-1$, $\Delta=k^2-x^2$ и $L=z+A(x)y$.
Тогда $L=ut+vs$, где $u=k-Ax$, $v=-x+Ak$.
Рассматривался анзац
\[
 H=P(x)L^2+Q(x)L+R(x)+M(x)y,
 \qquad H(C_3,C_8,C_7)=t,
\]
с $M=mx^e$. Сравнение коэффициентов в базисе $1,t,s$ даёт
необходимое тождество
\[
 v+P\bigl(u^3-3uv^2-xv^3\bigr)+mx^e\Delta=0.
\]
Его коэффициенты проверялись методом базисов Грёбнера. Для
$A=cx^j$, $j=0,1,2,3$, $1\leq e\leq16$, при неотрицательной
ожидаемой степени $P\leq6$ все исследованные идеалы равны $(1)$.
Отдельно проверены нечётные $A$ степеней не выше 1, 3, 5 и нечётные
$P$ степеней не выше 1, 3 (а также случаи $(1,5),(3,5)$):
во всех этих системах также получен базис $[1]$.
Степень $e$ во втором переборе равна $3\deg A+\deg P+3$.
Это точный отрицательный результат для выписанных конечных систем.
Он не исключает произвольных $A,P,M$, иных степеней, иных относительных
координат или конструкций с сохранением другой первой координаты.

\subsection{Полином параметра для $H(4,5,7)$}
Точная линейная алгебра по мономам показала, что минимальная общая
степень полинома $F$ с $F(C_4,C_5,C_7)=t$ равна 3, причём в степени
не выше 3 решение единственно:
\[
 F(x,y,z)=(x^2-2)z-y^3+3y.
\]
Этот конкретный $F$ не является координатой: его градиент обращается
в нуль при $x^2=2$, $y^2=1$, $z=0$. Следовательно, любое выпрямление
данной параметризации должно использовать полином параметра общей
степени как минимум 4 (или после аффинной перепараметризации ---
соответствующий аффинный образ полинома степени как минимум 4).
Это не препятствие существованию такого полинома большей степени.
